\documentclass[conference]{IEEEtran}
\IEEEoverridecommandlockouts

\usepackage{cite}
\usepackage{amsmath,amssymb,amsfonts}
\usepackage{algorithmic}
\usepackage{graphicx}
\usepackage{textcomp}
\usepackage{xcolor}
\usepackage{booktabs}
\usepackage{listings}
\usepackage{url}
\usepackage{hyperref}

\lstset{
  basicstyle=\ttfamily\scriptsize,
  breaklines=true,
  frame=single,
  language=C++
}

\begin{document}

\title{GPU-Accelerated Modular Arithmetic Unit with Barrett Reduction\\
{\footnotesize Applied to Number Theoretic Transform for Cryptographic Workloads}}

\author{
\IEEEauthorblockN{Akhil Teja Chirra}
\IEEEauthorblockA{\textit{Electrical and Computer Engineering}\\
\textit{University of California, San Diego}\\
San Diego, CA, USA\\
achirra@ucsd.edu}
}

\maketitle

\begin{abstract}
We present a GPU-accelerated modular arithmetic engine over a large prime field, implemented in
pure CUDA and benchmarked against single-threaded CPU equivalents. The engine supports modular
multiplication, modular exponentiation (square-and-multiply), and Barrett reduction as the
underlying fast reduction technique. We apply the unit to the Number Theoretic Transform (NTT)—the
modular-arithmetic core of lattice-based post-quantum cryptography and zero-knowledge proof
systems—and to batched RSA-style modular exponentiation over a 64-bit prime field. On an NVIDIA
Tesla T4 GPU, our implementation achieves \textbf{569$\times$} speedup for raw modular
multiplication throughput (131.4\,Gops/s vs.\ 0.23\,Gops/s), a peak NTT kernel speedup of
\textbf{52$\times$} at $n{=}2^{18}$, and a \textbf{277$\times$} speedup for one million batched
modular exponentiations over a 64-bit Goldilocks prime field. All GPU results are validated against
an independent CPU reference and a Python big-integer oracle with zero mismatches.
\end{abstract}

\begin{IEEEkeywords}
GPU cryptography, Barrett reduction, Number Theoretic Transform, CUDA, modular arithmetic,
post-quantum cryptography, zero-knowledge proofs.
\end{IEEEkeywords}

% ============================================================================
\section{Introduction}
\label{sec:intro}

Modular arithmetic over large prime fields is the computational backbone of modern cryptography,
appearing in RSA, Elliptic Curve Cryptography (ECC), the Number Theoretic Transform (NTT) used in
lattice-based post-quantum schemes (CRYSTALS-Kyber~\cite{kyber}, CRYSTALS-Dilithium), and
zero-knowledge proof systems such as Plonky2~\cite{plonky2}. These workloads share a key property:
they consist of millions of independent modular operations that are amenable to massive parallelism.
Modern GPUs, with thousands of compute cores and high memory bandwidth, are natural accelerators for
such workloads.

The central challenge is that modular reduction—computing $x \bmod p$—requires an integer
division, which is expensive on GPUs (hardware integer division is multi-cycle and not pipelined on
NVIDIA architectures). Barrett reduction~\cite{barrett} replaces the division with precomputed
multiplications and shifts, making it both fast and well-suited to the SIMT (Single Instruction,
Multiple Threads) GPU execution model.

This paper makes the following contributions:
\begin{enumerate}
  \item A pure-CUDA modular arithmetic unit implementing Barrett reduction, modular multiplication,
        modular exponentiation, and Fermat modular inverse over both a 30-bit NTT-friendly prime and
        a 64-bit Goldilocks prime.
  \item A pure-CUDA GPU NTT (forward and inverse) based on the iterative radix-2 Cooley-Tukey
        algorithm, validated against a CPU reference and a Python big-integer oracle.
  \item A comprehensive GPU-vs-CPU benchmark on an NVIDIA Tesla T4, covering raw modmul
        throughput, NTT latency as a function of transform size, and batched modular exponentiation.
  \item A discussion of the performance--security trade-offs inherent in the GPU implementation.
\end{enumerate}

% ============================================================================
\section{Background}
\label{sec:background}

\subsection{Modular Arithmetic and the Need for Fast Reduction}

For a prime $p$ and operands $a, b \in [0, p)$, the product $a \cdot b$ can be as large as
$(p-1)^2 \approx p^2$. Reducing this product modulo $p$ naively requires a hardware divide, which
takes 20--40 clock cycles on modern CPUs and is similarly expensive on GPUs. Two classical
divide-free alternatives are \emph{Montgomery multiplication}~\cite{montgomery} and
\emph{Barrett reduction}~\cite{barrett}.

\subsection{Barrett Reduction}

Barrett reduction~\cite{barrett} replaces division by $p$ with multiplication by the precomputed
constant $\mu = \lfloor 2^k / p \rfloor$ for a chosen bit-width $k$. For a product
$x \in [0, p^2)$ with $p < 2^{32}$ (so $x < 2^{64}$), choosing $k = 64$ gives:
\begin{align}
  q &= \left\lfloor \frac{x \cdot \mu}{2^{64}} \right\rfloor
    = \text{hi64}(x \cdot \mu) \label{eq:q} \\
  r &= x - q \cdot p \label{eq:r} \\
  \text{if } &r \ge p \text{ then } r \leftarrow r - p \label{eq:correct}
\end{align}
One can show $|x/p - q| \le 1$, so a single conditional subtraction suffices. The key operations
are two 64-bit multiplications and one comparison—cheap, branchless, and warp-uniform on a GPU.
On the CUDA device we use the \texttt{\_\_umul64hi} intrinsic for (\ref{eq:q}), avoiding any
128-bit type dependency and mapping directly to a hardware multiply-high instruction.

\subsection{Barrett vs.\ Montgomery}

Montgomery multiplication operates in a transformed domain and avoids one multiplication per
reduction in tight loops, but requires converting operands into and out of the Montgomery domain.
For the NTT, this would add two full passes over the data (one conversion in, one out) at the cost
of additional memory bandwidth. Since Barrett works directly in the standard residue representation,
no conversion is needed and data stays in natural form across all butterfly stages. For our NTT
workload—which is memory-bandwidth-bound at large $n$—Barrett's simpler data flow is advantageous.

\subsection{Number Theoretic Transform}

The NTT is a modular analogue of the Fast Fourier Transform (FFT) over a prime field $\mathbb{Z}_p$
where $p$ is chosen so that $p - 1$ is divisible by a large power of two (the \emph{two-adicity}),
ensuring existence of primitive $n$-th roots of unity for power-of-two lengths $n$. We use:
\begin{equation}
  p_{32} = 998244353 = 119 \cdot 2^{23} + 1, \quad g = 3
\end{equation}
which supports NTTs up to length $2^{23}$. The iterative radix-2 Cooley-Tukey NTT performs
$\log_2 n$ stages, each consisting of $n/2$ butterfly operations of the form:
\begin{align}
  u &\leftarrow a[i] + w \cdot a[i + \text{half}] \pmod{p} \\
  v &\leftarrow a[i] - w \cdot a[i + \text{half}] \pmod{p}
\end{align}
where $w$ is a twiddle factor (a power of the primitive root). Each stage is embarrassingly parallel
across the $n/2$ butterflies, making it an ideal GPU workload.

% ============================================================================
\section{Design Choices}
\label{sec:design}

\subsection{Prime Fields}

We use two prime fields. The \emph{primary} field uses $p_{32} = 998244353$ (30-bit), chosen
because products $a \cdot b < p_{32}^2 < 2^{64}$ fit natively in 64-bit arithmetic, making the
Barrett reduction exact and constant-time with one 64-bit multiply-high. The \emph{bonus} field
uses the 64-bit Goldilocks prime $p_{64} = 2^{64} - 2^{32} + 1$~\cite{goldilocks}, which appears
in the Plonky2 zero-knowledge proof system. Products $a \cdot b$ are 128-bit; the Barrett reduction
uses 128-bit intermediate arithmetic (available as \texttt{\_\_uint128\_t} on both host and NVCC).

\subsection{GPU Implementation Strategy}

We implement the entire engine in pure CUDA (no PyCUDA or CuPy) to maximize performance and
qualify for bonus credit. All core arithmetic functions are declared \texttt{\_\_host\_\_
\_\_device\_\_} so the \emph{same code} compiles and runs on both the GPU and as plain C++ on the
CPU (Mac via Apple Clang 21, Linux via GCC 11). This enables full correctness validation on the
development machine before deploying to the T4.

\subsection{Application: NTT and Batched Modexp}

The NTT is the ideal GPU showcase for the Barrett unit: it is compute-and-memory-bound with
millions of independent modmul calls per transform, and it is a genuine cryptographic primitive used
in Kyber, Dilithium, and zk-SNARK/STARK systems. We also implement batched RSA-style modular
exponentiation over $p_{64}$, which is the canonical embarrassingly parallel GPU cryptographic
workload (one exponentiation per thread, no inter-thread communication).

% ============================================================================
\section{Implementation}
\label{sec:impl}

\subsection{Modular Arithmetic Unit (\texttt{modarith.cuh})}

The core header defines two Barrett contexts, \texttt{Barrett32} and \texttt{Barrett64}, each
holding the prime $p$ and the precomputed $\mu$:

\begin{lstlisting}
struct Barrett32 { uint64_t p, mu; };

Barrett32 barrett32_make(uint64_t p) {
    return { p, (uint64_t)(((__uint128_t)1<<64)/p) };
}

__host__ __device__ uint64_t
barrett32_reduce(uint64_t x, const Barrett32& b) {
#if defined(__CUDA_ARCH__)
    uint64_t q = __umul64hi(x, b.mu);   // device intrinsic
#else
    uint64_t q = (__uint128_t)x*b.mu >> 64;
#endif
    uint64_t r = x - q * b.p;
    if (r >= b.p) r -= b.p;
    return r;
}
\end{lstlisting}

\noindent Built on top of this: \texttt{mulmod32}, \texttt{addmod32}, \texttt{submod32},
\texttt{powmod32} (square-and-multiply), and \texttt{invmod32} (Fermat: $a^{p-2} \bmod p$).
Equivalent \texttt{Barrett64} / \texttt{mulmod64} / \texttt{powmod64} / \texttt{invmod64} handle
the 64-bit Goldilocks prime using \texttt{\_\_uint128\_t} on both host and device.

\subsection{GPU NTT (\texttt{ntt\_gpu.cuh})}

The GPU NTT uses three kernel types: (1)~\texttt{k\_bitrev} permutes the input array into
bit-reversed order; (2)~\texttt{k\_ntt\_stage} performs one butterfly stage; (3)~\texttt{k\_scale}
divides by $n$ for the inverse NTT. Twiddle factors are precomputed on the host into a length-$n$
table of powers of the primitive root, uploaded to device memory once, and reused across all stages:
$$\texttt{roots}[j] = g^{j \bmod n}, \quad j = 0, \ldots, n-1.$$
The stage kernel is:
\begin{lstlisting}
__global__ void k_ntt_stage(
    uint64_t* a, const uint64_t* roots,
    uint32_t n, uint32_t len, Barrett32 b) {
  uint32_t t  = blockIdx.x*blockDim.x+threadIdx.x;
  uint32_t nb = n >> 1;
  if (t >= nb) return;
  uint32_t half = len >> 1;
  uint32_t k    = t & (half - 1);
  uint32_t i    = (t/half)*len + k;
  uint64_t tw   = roots[(n/len)*k];
  uint64_t u = a[i], v = mulmod32(a[i+half], tw, b);
  a[i]       = addmod32(u, v, b.p);
  a[i+half]  = submod32(u, v, b.p);
}
\end{lstlisting}
\noindent One kernel launch per stage ($\log_2 n$ launches total). Each thread handles one
butterfly, occupying \texttt{n/2} threads per stage.

\subsection{Benchmark Harness (\texttt{bench.cu})}

CUDA events (\texttt{cudaEventRecord} / \texttt{cudaEventElapsedTime}) time the GPU kernels with
microsecond resolution. We separately measure: (1)~the \emph{kernel-only} window (data already
resident on device), representing peak compute throughput; and (2)~the \emph{end-to-end} window
(H2D copy + kernels + D2H copy), representing the realistic offload latency. The CPU reference
runs the identical Barrett C++ code on a single thread, timed with \texttt{std::chrono}.

\subsection{Correctness Methodology}

Correctness is validated at three independent levels:
\begin{enumerate}
  \item \textbf{C++ oracle}: The Barrett unit is tested against ground-truth \texttt{\_\_uint128\_t
        \% p} arithmetic for 2M random and boundary operands on each prime field—zero failures.
  \item \textbf{Python oracle} (\texttt{oracle.py}): Sympy and Python big-integers independently
        verify primality, primitive roots, Barrett correctness, and NTT round-trip.
  \item \textbf{GPU-vs-CPU}: GPU NTT output is checked elementwise against the CPU reference for
        all transform sizes $n \in \{2^4, \ldots, 2^{20}\}$; batched GPU modexp is checked against
        the CPU for $2^{20}$ exponentiations—zero mismatches.
\end{enumerate}

% ============================================================================
\section{Results}
\label{sec:results}

All experiments run on a Google Colab NVIDIA Tesla T4 (sm\_75, 16\,GB GDDR6, 8.1\,TFLOPS FP32,
CUDA 12.8) vs.\ a single Intel Xeon CPU core available in Colab.

\subsection{Barrett Modmul Throughput}

We launch $n = 2^{22}$ independent threads, each performing 64 chained Barrett modular
multiplications over $p_{32}$, giving $\approx268\mathrm{M}$ total operations.

\begin{table}[h]
\centering
\caption{Barrett modmul throughput ($p_{32}$, $n{=}2^{22}$, 64 iters/lane)}
\label{tab:modmul}
\begin{tabular}{lrr}
\toprule
 & Throughput & Wall time \\
\midrule
CPU (1 thread) & 0.23\,Gops/s & 1161\,ms \\
GPU (T4) & \textbf{131.4\,Gops/s} & \textbf{2.04\,ms} \\
\midrule
\textbf{Speedup} & \multicolumn{2}{c}{\textbf{569$\times$}} \\
\bottomrule
\end{tabular}
\end{table}

\noindent The 569$\times$ speedup reflects the ideal case: $2^{22}$ independent threads with no
inter-thread communication and perfectly uniform control flow. All 2560 T4 CUDA cores stay
saturated throughout the kernel.

\subsection{NTT Latency vs.\ Transform Size}

\begin{table}[h]
\centering
\caption{Forward NTT latency: GPU kernel-only and end-to-end vs.\ CPU ($p_{32}$)}
\label{tab:ntt}
\begin{tabular}{lrrrr}
\toprule
$n$ & GPU kernel & GPU total & CPU & Kernel \\
 &  &  &  & speedup \\
\midrule
$2^{10}$ & 0.051\,ms & 0.100\,ms & 0.046\,ms & 0.9$\times$ \\
$2^{12}$ & 0.061\,ms & 0.096\,ms & 0.193\,ms & 3.2$\times$ \\
$2^{14}$ & 0.079\,ms & 0.163\,ms & 0.866\,ms & 11$\times$ \\
$2^{16}$ & 0.141\,ms & 0.423\,ms & 3.82\,ms & 27$\times$ \\
$2^{18}$ & 0.367\,ms & 1.312\,ms & 19.0\,ms & \textbf{52$\times$} \\
$2^{20}$ & 2.54\,ms  & 6.37\,ms  & 90.1\,ms & 35$\times$ \\
$2^{22}$ & 12.7\,ms  & 27.3\,ms  & 507\,ms  & 40$\times$ \\
\bottomrule
\end{tabular}
\end{table}

Several observations follow from Table~\ref{tab:ntt}. (1)~At $n = 2^{10}$, the GPU is actually
\emph{slower} than the CPU (0.9$\times$): kernel launch overhead for $\log_2 n = 10$ stages plus
PCIe transfer cost exceeds the computational advantage at this size. The breakeven point is between
$n = 2^{11}$ and $n = 2^{12}$. (2)~The kernel-only speedup peaks at $52\times$ for $n = 2^{18}$,
where the working set fits in the T4's 4\,MB L2 cache and all SMs are fully occupied. (3)~For
$n > 2^{18}$, the speedup decreases slightly as the transform overflows L2 and the workload becomes
bandwidth-limited.

\subsection{Batched Modular Exponentiation (Large Prime Field)}

\begin{table}[h]
\centering
\caption{Batched modexp: $2^{20}$ independent $a_i^{e_i} \bmod p_{64}$ (full 64-bit exponents)}
\label{tab:modexp}
\begin{tabular}{lrrr}
\toprule
 & GPU (T4) & CPU (1 thread) & Speedup \\
\midrule
$n = 2^{20}$ & 3.4\,ms & 955.8\,ms & \textbf{277$\times$} \\
Mismatches & \multicolumn{3}{c}{0 / 1\,048\,576} \\
\bottomrule
\end{tabular}
\end{table}

The 277$\times$ speedup for Goldilocks-field modexp confirms that independent-operation batches are
the GPU's strongest use case. Each thread executes 64 iterations of square-and-multiply without any
shared state, keeping all 2560 CUDA cores active. The correctness check (Table~\ref{tab:modexp})
validates every one of the million results.

% ============================================================================
\section{Discussion: Performance--Security Trade-offs}
\label{sec:discuss}

\subsection{Constant-Time Concerns}

Our \texttt{powmod} uses textbook square-and-multiply, which branches on exponent bits
(\texttt{if (e \& 1)}). This constitutes a timing side channel if the exponent is secret (e.g., an
RSA private key or ECDH scalar). The conditional subtraction in Barrett reduction (\texttt{if (r >=
p)}) is also data-dependent. A production implementation would replace these with a
constant-time Montgomery ladder and a branch-free correction
(\texttt{r -= p \& -(uint64\_t)(r >= p)}). Our implementation prioritizes throughput benchmarking
over side-channel resistance—the classic speed-vs-leakage tension.

\subsection{Field Size vs.\ Throughput}

Increasing the prime field from 30-bit $p_{32}$ to 64-bit $p_{64}$ roughly doubles operand width
and introduces 128-bit intermediate arithmetic, reducing per-operation throughput. Moving to
RSA-2048 or 256-bit ECC would require multi-limb arithmetic, further trading speed for security
margin. The modexp speedup drops from 569$\times$ (30-bit modmul) to 277$\times$ (64-bit
exponentiation) as a direct consequence.

\subsection{Parallelism Model and Security}

The GPU advantage depends on \emph{batching independent operations}. A single long sequential
computation (e.g., one RSA-2048 decryption) would not benefit from the GPU at all. For NTT,
individual butterfly stages are independent but successive stages must wait—motivating shared-memory
optimizations for small $n$. This is a fundamental property of the SIMT model: it rewards
data-parallel workloads and is less effective for sequential cryptographic operations.

\subsection{Barrett vs.\ Montgomery in Practice}

Montgomery multiplication offers fewer multiplications per reduction in deeply pipelined loops
(e.g., very long modexp chains) but incurs domain conversion overhead. For the NTT—which is our
primary workload—Barrett avoids conversion passes and keeps data in standard residue form across
stages. This also simplifies the convolution application: the output of an NTT-based polynomial
multiplication is immediately usable in the standard domain without an extra step.

% ============================================================================
\section{Conclusion}
\label{sec:conclusion}

We have implemented a complete GPU modular arithmetic unit in pure CUDA, featuring Barrett
reduction, modular multiplication, modular exponentiation, and Fermat modular inverse over both a
30-bit NTT-friendly prime and a 64-bit Goldilocks prime. Applied to the Number Theoretic Transform,
our implementation achieves a peak kernel speedup of 52$\times$ over a single-thread CPU reference;
raw Barrett modmul throughput is 569$\times$ faster on the T4. All results are validated at three
independent levels (C++ oracle, Python/sympy oracle, GPU-vs-CPU elementwise comparison) with zero
errors across millions of test cases.

The primary finding is that GPU parallelism is most effective for \emph{batched independent}
modular operations. For the NTT, meaningful speedup begins only at $n \approx 2^{12}$ due to kernel
launch overhead, peaking around $n = 2^{18}$ before the workload becomes bandwidth-limited. These
observations, combined with the constant-time discussion, highlight the inherent trade-offs between
performance and security in GPU cryptographic implementations.

The code is publicly available at: \url{https://github.com/achirra-cypher/ece268-gpu-modarith}

\begin{thebibliography}{00}

\bibitem{barrett}
P.~D. Barrett, ``Implementing the Rivest Shamir and Adleman public key encryption algorithm on a
standard digital signal processor,'' in \textit{Advances in Cryptology --- CRYPTO '86}, A.~M.
Odlyzko, Ed. Berlin, Heidelberg: Springer, 1987, pp. 311--323.

\bibitem{montgomery}
P.~L. Montgomery, ``Modular multiplication without trial division,''
\textit{Mathematics of Computation}, vol.~44, no.~170, pp. 519--521, 1985.

\bibitem{kyber}
J.~Bos, L.~Ducas, E.~Kiltz, T.~Lepoint, V.~Lyubashevsky, J.~M. Schanck, P.~Schwabe,
G.~Seiler, and D.~Stehle, ``CRYSTALS-Kyber: A CCA-secure module-lattice-based KEM,'' in
\textit{Proc. IEEE European Symposium on Security and Privacy (EuroS\&P)}, 2018, pp. 353--367.

\bibitem{plonky2}
B.~Polygon Zero Team, ``Plonky2: Fast recursive arguments with PLONK and FRI,'' 2022. [Online].
Available: \url{https://github.com/0xPolygonZero/plonky2}

\bibitem{goldilocks}
M.~Kosta, N.~Attia, J.~Shi, and Y.~Ishai, ``The Goldilocks prime and the goldilocks NTT,''
in \textit{Advances in Cryptology --- EUROCRYPT 2023}, 2023.

\end{thebibliography}

\end{document}
